\chapter{Modules}
\label{ch:modules}

A module is a "container" for code; it shares the same behavior of
scopes~(\see{ch:scopes}): it can hold code, variables, functions, constants,
types definition etcetera; in fact it can be seen as a "named scope".

Each file is implicitly a module, whose name is the filename containing it.

\begin{minted}{\giovilexer}
# my_project/src/Foo.giovi
"hello from Foo module" -> _out
const SOME_CONSTANT Real <- 3.14
someFunction() => Int {
	"hello from Foo.someFunction" -> _out
	=> 69
}
\end{minted}

\section{Inner modules}

A inner sub-module can be declared with the \code{module} keyword.

\begin{minted}{\giovilexer}
# my_project/src/Foo.giovi
"hello from Foo module" -> _out
module Bar {
	"hello from Foo.Bar module" -> _out
	module Baz {
		"hello from Foo.Bar.Baz module" -> _out
	}
}
\end{minted}

\nb{The body of an inner module is not executed until it is imported! See
\See{sec:import}.}

\section{Importing modules}
\label{sec:import}

The \code{import} statement is used to "import" one or more modules.

\begin{minted}{\giovilexer}
module Foo {
	"hello from Foo module" -> _out
}
import Foo # "hello from Foo module"
\end{minted}

\subsection{Import entry}

An import entry can be a qualified identifier (for native and inner modules) or
a file URI for file modules. To later refer to a file module, it must be
aliased first.

\begin{minted}{\giovilexer}
import [
	giovi.util.Math,
	'file://my_project/my_module.giovi' as MyModule
]
giovi.util.Math -> _out
MyModule -> _out
\end{minted}

Import entries URIs are written using the text literals syntax
(\see{sec:text_syntax}) and, as such, it allows for text interpolation, which
is advised to be used carefully.

\begin{minted}{\giovilexer}
nameOfMyModule Text <- "my_other_module"
import "file://my_project/${nameOfMyModule}.giovi"
\end{minted}

It's possible to import to a inner module of a file module with a single import
statement; however the URI can only appear as the first entry

\begin{minted}{\giovilexer}
import 'file://my_project/my_module.giovi'.MySubModule.MySubSubModule
\end{minted}

\subsection{Import behaviour}

The \code{import} keyword may be deceiving as the import statement doesn't
really import a module like in many other languages, but rather it executes its
body and then gives you access to it and its declarations.

\nb{The import of the same module can be performed multiple times, meaning that
its body will be executed for each import.}

\todo{menzionare funzione util per caching moduli}

\begin{minted}{\giovilexer}
module MyMath {
	"hello from MyMath module" -> _out
	const PI <- 3.14
}
import MyMath # "hello from MyMath module"
MyMath.PI -> _out # 3.14
\end{minted}

\paragraph{Lazy imports} A peculiarity of the import statement is that it can
be performed anywhere like any other statement; the lifetime of the imported
module alias, like all other declarations (e.g. variables, constants,
functions), will be bound to the scope within which the import is executed.\\*
This allows the flexibility of deferring the import to whenever it is actually
needed.

\begin{minted}{\giovilexer}
import giovi.util.Random as Random # immediately imported

if Random.real() > 0.69 do {
	import giovi.util.Math as Math
	Math.PI -> _out # 3.14
}
\end{minted}

